\section{相关工作}
\label{sec:ink-related-work}

\subsection{模型合并}

模型合并旨在将多个由同一预训练模型独立微调得到的专家模型整合为单一多任务模型，从而避免联合训练以及多模型推理带来的额外开销。早期方法主要在参数空间中进行组合，例如模型权重平均~\citep{wortsman2022model} 和任务向量算术~\citep{ilharco2023editing} 通过聚合专家参数更新实现知识融合。随后，研究逐渐关注不同任务更新之间的干扰问题，并从参数稀疏化、符号冲突以及任务子空间等角度提升合并效果。TIES-Merging 通过裁剪冗余更新并协调参数符号缓解任务间冲突~\citep{yadav2023ties}；TSV-M、Iso-C/CTS、WUDI-Merging 和 Essential Subspace Merging 等方法进一步利用任务更新的谱结构、共享子空间或任务特定子空间改善能力保留与组合~\citep{gargiulo2024task,marczak2025no,cheng2025whoever,li2026essential}。

除参数空间方法外，另一类工作利用额外信息自适应地调整合并过程。例如，AdaMerging 通过无标签数据上的熵优化学习任务级或层级合并系数~\citep{yang2024adamerging}。

\subsection{几何感知模型合并}

近年来，几何感知模型合并方法进一步将合并过程从简单参数距离扩展到任务相关几何空间，通过非各向同性度量区分不同参数方向的重要性。Fisher Merging 利用专家模型附近的 Fisher 信息刻画参数扰动对任务性能的影响~\citep{matena2022fisher}；RegMean 根据任务数据中的输入激活统计匹配不同模型的层输出，从而保留任务中频繁访问的输入方向~\citep{jin2023regmean}。MaTS 则将上述方法统一视为任务参数子空间中的模型匹配问题，并利用 Kronecker 分解和共轭梯度求解包含输入侧与输出侧结构的双侧线性系统~\citep{tam2024mats}。Chain of Merges 进一步考虑逐层合并过程中前序层变化导致的输入分布漂移，通过动态更新统计信息改善连续层之间的匹配~\citep{buzzega2025chain}。

这些方法的共同目标是通过任务相关几何判断哪些方向需要在合并过程中被保护。然而，它们所使用的几何通常来源于专家模型自身的参数统计、局部敏感性或任务数据分布。例如，RegMean 描述任务输入方向的覆盖程度，Fisher 和 MaTS 则描述专家附近参数扰动对任务目标的敏感程度。换言之，这些方法主要刻画专家模型中哪些方向具有重要性，但并不显式描述在当前合并模型下，为匹配某个专家仍需要修正哪些功能方向。

本文进一步将几何对象从固定的专家重要性几何扩展为动态的专家--当前模型修正几何。具体而言，\method{} 根据专家模型与当前合并模型之间的行为差异构造动态功能修正，并使该几何随着合并状态变化而重新估计。该方法从保护专家已有能力的重要方向，转向刻画当前模型为了接近专家仍需发生修正的功能方向。

\paragraph{与 MaTS 的详细比较。}
MaTS 用半正定矩阵 $C_t$ 表示专家 $t$ 的任务参数子空间，将合并写成
\begin{equation}
  \mathcal J_C(\theta)
  =
  \frac12\sum_{t=1}^T
  (\theta-\theta_t)^\top C_t(\theta-\theta_t),
  \qquad
  \left(\sum_tC_t\right)\theta^\star
  =
  \sum_tC_t\theta_t.
  \label{eq:ink-mats-subspace}
\end{equation}
对于线性层，K-FAC 近似令
\begin{equation}
  C_t^\ell\approx G_{t,F}^\ell\otimes A_t^\ell,
  \qquad
  G_{t,F}^\ell
  =
  \mathbb E\left[\delta_t^\ell(\delta_t^\ell)^\top\right],
\end{equation}
其中 $\delta_t^\ell$ 是任务损失对层输出的梯度。对应的双侧系统为
\begin{equation}
  \sum_tA_t^\ell X^\star G_{t,F}^\ell
  =
  \sum_tA_t^\ell X_t^\ell G_{t,F}^\ell.
  \label{eq:ink-mats-bilateral}
\end{equation}
Kronecker 乘积之和通常不能化为因子和的单个 Kronecker 乘积，因此 MaTS 将左侧视为线性算子，并用共轭梯度避免显式求逆。

\method{} 与 MaTS/K-FAC 在代数形式和求解器上最接近，但算子系数的语义与生命周期不同。MaTS 的 $G_{t,F}^\ell$ 是专家附近任务损失的输出梯度协方差，用固定任务子空间衡量哪些扰动敏感；\method{} 的 $G_{t,\mathrm{ink}}^{\ell,(r)}$ 则由专家--当前模型 KL 残差诱导，衡量当前模型尚需沿哪些功能方向修正，并在每次接受更新后改变。因此，我们复用的是双侧线性系统与共轭梯度这一计算骨架，而不是对同一个固定 Fisher 系统做更精确求解。
